\subsection{Dedekind overrings and all models of certain strictly local fields}

\begin{lemma}[Proved; the generic upper bound passes to Dedekind overrings]
\label{A:Dedekind-overring-cd}
Let $D$ be a Dedekind domain, let $p\in D^\times$ be prime,
and let $K=\Frac D$ have finite $p$-cohomological dimension $n$.
Suppose $\cd_p(\Spec D)\leq n$. Then every ring
$D\subset E\subset K$ satisfies
\[
 \cd_p(\Spec E)\leq n.
\]
Every such $E$ is a Dedekind domain or a field. It need not be a
localization of $D$ at a multiplicative subset.
\end{lemma}
\begin{proof}
If $D$ is a field there is nothing to prove. Otherwise $n\geq1$.
Krull--Akizuki, in the overring form of
\citet[Tag 00PG, Lemma 10.119.12]{A_Stacks}, makes $E$ noetherian.
We recall explicitly the prime correspondence that we need.
For a nonzero prime $\mathfrak q\subset E$, its contraction
$\mathfrak p=\mathfrak q\cap D$ is nonzero: writing a nonzero
$x\in\mathfrak q$ as $a/b$ with nonzero $a,b\in D$ gives
$a=bx\in\mathfrak p$. The inclusion
$D_{\mathfrak p}\subset E_{\mathfrak q}\subset K$ is an inclusion
of local overrings of the discrete valuation ring $D_{\mathfrak p}$,
and its uniformizer remains a nonunit. The only proper overring
of a discrete valuation ring in its fraction field is the field.
Consequently
\[
 E_{\mathfrak q}=D_{\mathfrak p}.
\]
This proves that $E$ is Dedekind when it is not a field. It also
shows that contraction is injective on nonzero primes, with
unchanged local discrete valuation rings and residue fields:
$\mathfrak q$ is recovered as the inverse image of the maximal
ideal of this common local ring. This is the overring correspondence
also stated in \citet[Lemma 1, p.~216]{A_Clark2009}.

Fix a finite separable extension $L/K$ containing $\mu_p$.
Let $D_L$ and $E_L$ be the integral closures of $D$ and $E$ in $L$.
They are Dedekind or fields and are finite over their respective
bases, by the trace argument in
Lemma~\ref{A:Dedekind-Milnor-obstruction}.
Since every element integral over $D$ is integral over $E$,
$D_L\subset E_L\subset L$. The preceding local-ring argument
identifies the nonzero primes of $E_L$ with a subset of those of
$D_L$, preserving the associated normalized discrete valuations.
The residue map for $E_L$ is therefore the composite
\[
 K_n^M(L)/p\longrightarrow
 \bigoplus_{\mathfrak p\in(\Spec D_L)^{(1)}}
 K_{n-1}^M(\kappa(\mathfrak p))/p
 \longrightarrow
 \bigoplus_{\mathfrak q\in(\Spec E_L)^{(1)}}
 K_{n-1}^M(\kappa(\mathfrak q))/p,
\]
where the second arrow is the projection to that subset.
The first arrow is surjective by
Theorem~\ref{A:Dedekind-Milnor-criterion} and the hypothesis on $D$;
hence the composite is surjective. The same criterion applied to
$E$ gives the asserted bound. If $E=K$, the conclusion is immediate.
No finite-type hypothesis on $D\to E$ has entered the argument.
\end{proof}

\begin{lemma}[Proved; strictly local rings lie in every noetherian model normalization]
\label{A:strict-normalization-rigidity}
Let $B$ be a strictly henselian noetherian local domain with a prime
$p$ invertible in $B$, and put $K=\Frac B$. If $R\subset K$ is any
noetherian domain with fraction field $K$, its integral closure
$\overline R$ in $K$ contains $B$. The normalization need not be
finite or noetherian.
\end{lemma}
\begin{proof}
Every $u\in B^\times$ has a $p^r$th root in $B^\times$ for each
$r\geq1$: the polynomial $X^{p^r}-u$ has a simple root over the
separably closed residue field and strict henselianity lifts it.
Thus every discrete valuation $v:K^\times\to\Z$ vanishes on
$B^\times$, since $v(u)$ is divisible by every $p^r$.
For $b$ in the maximal ideal of $B$, the element $1+b$ is a unit.
If $v(b)<0$, the valuation inequality with unequal values gives
$v(1+b)=v(b)<0$, a contradiction. Therefore $B$ lies in every
discrete valuation ring of $K$.

The Mori--Nagata theorem
\citep[Theorem 4.10.5, pp.~90--92]{A_SwansonHuneke2006}
states that the integral closure of a noetherian domain in its fraction field is a Krull ring. In particular,
$\overline R$ is the intersection in $K$ of its height-one
discrete valuation local rings. The preceding paragraph puts $B$
in each of these rings, and hence in $\overline R$. If
$\overline R=K$, the same containment is immediate. We have not
replaced a possibly infinite normalization by a finite one. The same global formulation is recorded in 
\citet[Chapitre 0, Remarque 23.2.8(iii), p.~220]{A_EGAIV1}.
\end{proof}

\begin{theorem}[Proved; A.2 for every model of a specified strictly local surface field]
\label{A:A2-all-surface-models}
Let $B$ be a strictly henselian noetherian local domain of dimension
two, let $p$ be invertible in $B$, and suppose that some nonzero
$t$ in its maximal ideal has quasi-excellent quotient $B/tB$.
Put $K=\Frac B$. Then A.2 holds for every noetherian domain
$R\subset K$ with $\Frac R=K$ and $p\in R^\times$.
The domain $R$ need not be normal, excellent, local, or of finite
type over a subring of $B$.

In particular, this covers every noetherian model of the fraction
field of every regular strictly henselian local surface of
equicharacteristic zero, including the Nagata--Gabber example.
\end{theorem}
\begin{proof}
We know $\cd_p(K)=2$ by Theorem~\ref{A:strict-field}.
The ring $B[1/t]$ has dimension at most one, and
Proposition~\ref{A:principal-completion-bound} gives
$\cd_p(B[1/t])\leq2$. Let $D$ be its integral closure in $K$.
Krull--Akizuki makes $D$ Dedekind or a field. Exact integral
pushforward on the small \'etale site gives
\[
 \cd_p(\Spec D)\leq2.
\]
Indeed, for the integral map $\nu:\Spec D\to\Spec B[1/t]$
and every $p$-primary torsion sheaf $\mathcal E$ on $\Spec D$,
$H^q(\Spec D,\mathcal E)=
H^q(\Spec B[1/t],\nu_*\mathcal E)$.

Let $\overline R$ be the integral closure of $R$ in $K$.
Lemma~\ref{A:strict-normalization-rigidity} gives
$B\subset\overline R$. Since $t\in K^\times$, choose
$0\neq a\in R$ so that both $t$ and $t^{-1}$ belong to $R[1/a]$.
For example, if $t=r/s$ with nonzero $r,s\in R$, we may initially
take $a=rs$. When $p=2$, strict henselianity puts an element
$i$ with $i^2=-1$ in $B$; enlarge $a$ by the denominator of an
expression for $i$ in $K=\Frac R$. Thus all residue fields of
$R[1/a]$ are nonorderable in this case.

Write $R_0=R[1/a]$ and $E=\overline R[1/a]$.
Integral closure commutes with localization, so $E$ is the
integral closure of $R_0$. It is integrally closed in $K$ and
contains $B[1/t]$, and therefore contains $D$.
Lemma~\ref{A:Dedekind-overring-cd} gives that $E$ is Dedekind
or a field and that $\cd_p(E)\leq2$.
The integral surjection $\Spec E\to\Spec R_0$ preserves
Krull dimension, so $\dim R_0\leq1$.

More generally, let $W\subset\Spec R_0$ be a nonempty affine
open and put $S=\Gamma(W,\mathcal O_W)\subset K$.
The normal ring $S^\nu$, the integral closure of $S$ in $K$,
contains $E$: every element of $E$ is integral over $R_0\subset S$.
Thus it is again an overring of $D$ and has cohomological
dimension at most two. If $S$ is a field then $S=K$ and the same
bound holds. Otherwise $S$ is a noetherian one-dimensional domain.
The sheaf comparison proved in
Theorem~\ref{A:curve-normalization} gives $\cd_p(S)\leq2$:
the unit into integral-normalization pushforward has a generically
zero cokernel whose cohomology vanishes above degree one.
The real-residue exception has already been excluded when $p=2$.
Hence the upper bound two holds on every such affine $W$.

Finally Proposition~\ref{A:spread} supplies a nonempty open on
which every nonempty affine subopen has cohomological dimension
at least $\cd_p(K)=2$. Intersect it with $\Spec R_0$.
Every nonempty open of $\Spec R$ contains a nonempty affine open
inside this intersection, and every further nonempty affine
subopen has cohomological dimension exactly two. This proves the
full quantifier in A.2.

For the last assertion, a regular parameter $t$ in a regular
strictly henselian local surface of equicharacteristic zero has
characteristic-zero discrete valuation quotient, hence excellent,
as checked in Theorem~\ref{A:A2-local-surface}.
\end{proof}

\begin{remark}[Search and scope]
We checked the full Mori--Nagata proof in Huneke--Swanson,
the Stacks Project's precise Krull--Akizuki overring statement,
and Clark's explicit Dedekind prime correspondence. The
cohomological comparison uses the residue criterion already proved
above. We make no priority claim for the resulting model theorem.
It removes all noetherian models of these specified surface fields
from the A.2 gate; it does not settle arbitrary one-dimensional
arithmetic fields or arbitrary strictly local fields in higher
dimension.
\end{remark}
